% TOP_PROBLEM: {"problemId": "problem.auslander-conjecture-on-affine-crystallographic-groups", "canonicalTitle": "Auslander Conjecture on Affine Crystallographic Groups", "releaseRank": 472, "primaryDomain": "geometry_topology", "primaryDomainLabel": "Geometry & topology", "displayStatus": "open_with_solved_subcases", "releaseStatus": "open_with_solved_subcases"}
% !TeX program = lualatex
% ATTEMPT_STATUS: unresolved
% Model: GPT-6 Astra Ultra; continuation dated 27 September 2026.
\documentclass[11pt]{article}
\usepackage[margin=25mm]{geometry}
\usepackage{fontspec}
\setmainfont{Latin Modern Roman}
\usepackage{amsmath,amssymb,mathtools}
\usepackage{unicode-math}
\setmathfont{Latin Modern Math}
\usepackage{xurl,hyperref}
\hypersetup{colorlinks=true,urlcolor=blue}
\setcounter{secnumdepth}{0}
\setlength{\parindent}{0pt}
\setlength{\parskip}{5pt}
\setlength{\emergencystretch}{3em}
\begin{document}
\section*{\texorpdfstring{Auslander Conjecture on Affine Crystallographic Groups}{Auslander Conjecture on Affine Crystallographic Groups}}\label{472}
\subsection{Definitions and mathematical statement}
An affine transformation of ${\mathbb{R}}^n$ has the form $x\mapsto Ax+b$ with $A\in\operatorname{GL}_n({\mathbb{R}})$. An affine crystallographic group is a group of such transformations acting properly discontinuously and cocompactly on ${\mathbb{R}}^n$. Virtually solvable means having a solvable subgroup of finite index. Auslander's conjecture asserts that every affine crystallographic group is virtually solvable.

In the manifold formulation, a complete affine manifold is a quotient by a free properly discontinuous affine action; completeness identifies its developing space with all of ${\mathbb{R}}^n$. For compact such manifolds, the equivalent target in the record is that the fundamental group be virtually polycyclic, meaning that a finite-index subgroup admits a finite subnormal series with cyclic factors. No preserved positive-definite metric is assumed: restricting to Euclidean isometries would discard the essential affine difficulty.
\par\smallskip\textit{Formulation and concept references:} \href{https://www.math.umd.edu/~wmg/gstom.pdf}{[S1]}.

\subsection{Short English statement}
Are groups acting properly and cocompactly by affine transformations always virtually solvable, with the corresponding compact complete manifold groups virtually polycyclic?
\subsection{Research attempt}

\paragraph{Scope and established results (27 September 2026).}
Both properness and cocompactness are essential hypotheses. Here
properness means that for every compact $K\subset\mathbb R^n$ only
finitely many $\gamma$ satisfy $\gamma K\cap K\ne\varnothing$.
This permits finite stabilizers in the group formulation. The manifold
formulation additionally requires a free action. We never infer
freeness from properness without a torsion-free hypothesis.

Goldman's text [S1] distinguishes the compact complete affine problem
from the existence of complete noncompact manifolds with nonsolvable
fundamental groups. Abels--Margulis--Soifer [E1] prove the conjecture
for $n<7$; their result is much stronger than the elementary subcases
proved here. This audit does not claim that a proof in unrestricted
dimension has been found. The arguments below explore the translation
subgroup, triangular linear holonomy, and compact quotient topology.

\paragraph{Linear and translational parts.}
Write an affine transformation as $(A,b)$ with action $x\mapsto Ax+b$.
Multiplication and inversion are
\[
 (A,b)(B,c)=(AB,b+Ac),\qquad
 (A,b)^{-1}=(A^{-1},-A^{-1}b).
\]
The linear-part map $\ell:\Gamma\to\operatorname{GL}_n(\mathbb R)$
is a homomorphism. Its kernel $T$ consists of translations and is
abelian. Identifying $T$ with its translation vectors gives
\[
 1\longrightarrow T\longrightarrow\Gamma
 \stackrel{\ell}{\longrightarrow}\ell(\Gamma)\longrightarrow1.
\]
Conjugation sends translation by $t$ to translation by $At$, so $T$
is normal and its real span $V$ is $\ell(\Gamma)$-invariant.

Properness implies that $T$ is a discrete subgroup of $\mathbb R^n$:
infinitely many distinct bounded translation vectors would violate
properness for a sufficiently large closed ball. A discrete subgroup
of a finite-dimensional real vector space is a lattice in its span,
and hence isomorphic to $\mathbb Z^r$ for some $r\leq n$.
One elementary justification for finite generation is to choose a real
basis $t_1,\ldots,t_r$ from $T$. Subtract integral combinations of
these vectors to move every element of $T$ into their bounded
parallelepiped. Its intersection with $T$ is finite by discreteness.
These finitely many representatives, together with the $t_i$, generate
$T$. A finitely generated torsion-free abelian group is free, and its
discreteness forces its rank to equal the dimension of its real span.

\paragraph{A full-rank translation subgroup settles a genuine subcase.}
\textbf{Proposition 472.1.} If $\Gamma$ acts properly on $\mathbb R^n$
and its translation subgroup has rank $n$, then $T$ has finite index
in $\Gamma$. In particular $\Gamma$ is virtually abelian and hence
virtually polycyclic.

\emph{Proof.}
Take a compact parallelepiped $Q$ for the full lattice $T$. For every
$\gamma\in\Gamma$ choose $t\in T$ so that $t\gamma(0)\in Q$.
Thus each left coset $T\gamma$ has a representative sending $0$ into
$Q$. Only finitely many such representatives exist: take a compact
set $K$ containing $Q\cup\{0\}$ and apply properness to it.
Therefore there are finitely many cosets.\hfill$\square$

Notice that cocompactness was not required once a full translation
lattice was given. The missing step in applying this proposition to
the general problem is not its proof; it is the false hope that every
affine crystallographic group must have a full translation lattice.
An explicit compact affine example below disproves that hope.

There is also a useful converse subcase. If $\ell(\Gamma)$ is finite,
then $T$ has finite index. Cocompactness for $\Gamma$ gives cocompactness
for $T$: enlarge a compact covering set by finitely many coset
representatives. A translation lattice of rank less than $n$ leaves
an unbounded quotient direction, as the projection to $\mathbb R^n/V$
is constant on its orbits. Thus finite linear holonomy forces
rank $T=n$ and again yields the conclusion.

\paragraph{Dimension one, including finite stabilizers.}
\textbf{Proposition 472.2.} Every affine crystallographic group of the
line is virtually infinite cyclic.

\emph{Proof.}
Let $\Gamma^+$ be the subgroup with positive linear part. It has index
at most two. An element $x\mapsto ax+b$ with $a>0$ and $a\ne1$ has
the fixed point $b/(1-a)$ and infinite order, so its powers give an
infinite stabilizer. This contradicts properness. Thus every element
of $\Gamma^+$ is a translation. The discrete translation subgroup of
the line is either zero or $c\mathbb Z$ for some $c>0$. The zero
case would make $\Gamma$ finite, and a finite group cannot cover the
unbounded line by translates of a compact set. Hence the second case
occurs.\hfill$\square$

The same reasoning must not be copied to higher dimension by replacing
$a\ne1$ with $A\ne I$: a nonidentity matrix can have eigenvalue one,
and an affine map with that linear part can be fixed-point-free.

\paragraph{A solvable linear-holonomy criterion.}
\textbf{Proposition 472.3.} If the linear image $\ell(\Gamma)$ is
virtually solvable, then $\Gamma$ is virtually solvable. In particular
this holds if a finite-index subgroup of the linear image is
simultaneously upper triangular over $\mathbb C$.

\emph{Proof.}
Take the inverse image $\Gamma_0$ of a solvable subgroup of finite
index in $\ell(\Gamma)$. If its image has derived length $d$, then
$\Gamma_0^{(d)}\subseteq T$. Since $T$ is abelian,
$\Gamma_0^{(d+1)}=1$. This proves the first assertion.

For the second, the diagonal map on the upper triangular group has
abelian image and kernel the upper unitriangular group $U$. Let $U_j$
be the subgroup whose entries on the first $j-1$ superdiagonals vanish.
Matrix multiplication gives $[U_i,U_j]\subseteq U_{i+j}$, and
$U_n=1$. Consequently repeated commutators eventually vanish in $U$.
The triangular group is an extension of this nilpotent group by an
abelian group, hence is solvable.\hfill$\square$

This establishes virtual solvability directly. It does not by itself
prove the virtually polycyclic conclusion for every finitely generated
solvable linear group: that general assertion is false. For example,
the group generated on the line by $x\mapsto2x$ and $x\mapsto x+1$
contains translations by every dyadic rational. It is a solvable affine
group, but those arbitrarily small nonzero translations violate
properness. Its translation subgroup $\mathbb Z[1/2]$ is not finitely
generated, so the group is not polycyclic. Proper crystallographic
actions supply additional restrictions in the established equivalence
discussed in [S1]; we do not replace those restrictions by solvability
alone.

\paragraph{A compact non-Euclidean example with too few translations.}
For integers $a,b,c$, define
\[
 g_{a,b,c}(x,y,z)=(x+a,\ y+b,\ z+c+ay).
\]
Direct substitution yields
\[
 g_{a,b,c}g_{a',b',c'}
 =g_{a+a',\,b+b',\,c+c'+ab'},\qquad
 g_{a,b,c}^{-1}=g_{-a,-b,-c+ab}.
\]
Thus these maps form the integral Heisenberg group, here given by an
explicit affine action on all of $\mathbb R^3$.

\textbf{Proposition 472.4.} This action is free, proper and cocompact.
The group is polycyclic and nilpotent of class two, is not virtually
abelian, and its subgroup of translations has rank two rather than three.

\emph{Proof.}
A fixed point would force $a=b=0$ from the first two coordinates and
then $c=0$ from the third. Hence the action is free.

If a compact set is contained in $[-M,M]^3$ and its image by
$g_{a,b,c}$ intersects it, the first two coordinates give
$|a|,|b|\leq2M$. The third gives
$|c|\leq2M+|a|M$. There are only finitely many integral triples
satisfying these bounds, proving properness.

Given any $(x,y,z)$, choose integers $a,b$ so that $x+a,y+b\in[0,1]$,
and then choose $c$ so that $z+c+ay\in[0,1]$. Its orbit meets the
unit cube, proving cocompactness.

The map $g_{a,b,c}\mapsto(a,b)$ has central kernel
$\{g_{0,0,c}\}$ and quotient $\mathbb Z^2$. More explicitly,
\[
 [g_{a,b,c},g_{a',b',c'}]=g_{0,0,ab'-a'b}.
\]
This proves nilpotence of class two. The subgroups with $a=0$ and then
$a=b=0$ give a subnormal series with three infinite cyclic factors,
proving polycyclicity. Any finite-index subgroup contains positive
powers of $g_{1,0,0}$ and of $g_{0,1,0}$. Their commutator is
$g_{0,0,mn}\ne1$, so that subgroup is not abelian.

Finally the linear matrix is
\[
 \begin{pmatrix}1&0&0\\0&1&0\\0&a&1\end{pmatrix}.
\]
It equals the identity exactly when $a=0$. The translation subgroup
therefore consists of vectors $(0,b,c)$ and has rank two.
\hfill$\square$

This example satisfies the conjecture, but defeats the stronger proposed
intermediate conclusions ``finite linear holonomy'' and ``virtually
abelian''. It also prevents an argument that averages a positive-definite
inner product over the linear holonomy without proving that averaging
is legitimate. A nontrivial unipotent matrix cannot preserve a
positive-definite form: a matrix preserving such a form is conjugate
to an orthogonal matrix and is diagonalizable over $\mathbb C$,
whereas the displayed nonidentity shear is not.

\paragraph{What fixed-point tests really imply.}
For $\gamma=(A,b)$, the fixed-point equation is
\[
 (I-A)x=b.
\]
If $1$ is not an eigenvalue of $A$, the equation has a unique solution.
Thus every infinite-order element of a proper affine group must have
$1$ as an eigenvalue of its linear part. In a free action this holds
for every nonidentity element, whether or not its order was specified.

Even the converse fixed-point test needs care: if $1$ is an eigenvalue,
there is a fixed point exactly when $b\in\operatorname{im}(I-A)$.
Writing $\lambda$ for a covector fixed by $A^T$, the necessary
condition $\lambda(b)=0$ for all such $\lambda$ is also sufficient,
by elementary duality between image and annihilator. Translation in a
direction detected by one of these covectors can remove all fixed points.
Freeness still does not imply properness of the whole group.

The spectral restriction alone does not force solvability. Consider the
action of $\operatorname{SL}_2(\mathbb R)$ on the three-dimensional
space of trace-zero matrices by $X\mapsto AXA^{-1}$. For every
$A$, the matrix $A-\tfrac12\operatorname{tr}(A)I$ is a fixed trace-zero
matrix when it is nonzero; if it vanishes, $A=\pm I$ and the whole
space is fixed. Thus every matrix in this representation has eigenvalue
one. Nevertheless its image contains a nonabelian free group.

Here is an explicit verification of the last assertion. The matrices
\[
 U=\begin{pmatrix}1&2\\0&1\end{pmatrix},\qquad
 V=\begin{pmatrix}1&0\\2&1\end{pmatrix}
\]
act on the real projective line by $t\mapsto t+2$ and
$t\mapsto t/(2t+1)$. Put $X=\{|t|>1\}\cup\{\infty\}$ and
$Y=\{|t|<1\}$. For every nonzero integer $m$,
$U^m(Y)\subset X$ and $V^m(X)\subset Y$: for the latter,
$|2mt+1|\geq2|t|-1>|t|$ when $|t|>1$, and the assertion at
infinity follows from $V^m(\infty)=1/(2m)$.
The ping-pong argument now shows that $U,V$ freely generate. To spell
out the usual endpoint issue, a nontrivial reduced word may be conjugated
to a pure power or to an alternating word beginning and ending in
nonzero powers of the same generator; in the latter case it carries
the other generator's domain strictly into that generator's domain.
It is not the identity. Pure powers are visibly nonidentity as well.
The projective action also excludes $-I$ from this free subgroup.
The adjoint representation, whose kernel is $\{I,-I\}$, therefore
remains faithful on it.

This is only a counterexample to a proposed algebraic implication from
the eigenvalue condition. Its linear action fixes the origin and is
not proper. It is not a counterexample to Auslander's conjecture.

\paragraph{A compactness obstruction from cohomological dimension.}
Suppose now that the action is free, proper and cocompact. The quotient
$M=\Gamma\backslash\mathbb R^n$ is a compact manifold without boundary
with contractible universal cover. Its lifted cellular chain complex
gives a free resolution of the trivial module over
$\mathbb F_2[\Gamma]$ of length $n$, using a finite-dimensional
triangulation of the smooth compact manifold. On the other hand,
\[
 H^n(\Gamma;\mathbb F_2)=H^n(M;\mathbb F_2)\ne0
\]
by the mod-two fundamental class. Hence
$\operatorname{cd}_{\mathbb F_2}\Gamma=n$.

A nontrivial free group has cohomological dimension one over
$\mathbb F_2$: its tree gives the upper bound, and its nonzero first
cohomology gives the lower bound. Therefore a nonabelian free group
cannot act freely, properly and cocompactly on $\mathbb R^n$ for $n>1$.
This explains why a proper free affine action in dimension three,
of the kind treated in the established Margulis-spacetime theory,
does not refute the compact conjecture. It does not exclude a much
larger group containing a free subgroup: the cohomological dimension
of a subgroup need not equal that of the ambient group.

\paragraph{Exact diagnostics and remaining gap.}
The accompanying integer-arithmetic script checks the Heisenberg
composition, inverse and commutator identities on a finite grid and
checks the adjoint matrices of the two displayed generators. These
checks detect transcription and convention errors; the universally
quantified proofs are the symbolic calculations above. No finite
matrix search certifies properness for an arbitrary infinite affine group.

The proved subcases are the full translation lattice case, finite
linear holonomy, dimension one, and virtually solvable (in particular
virtually triangular) linear holonomy. The explicit compact example
shows why these do not exhaust the problem. The first unresolved step
is to exclude a nonsolvable linear image using both the translational
part and cocompactness; neither the eigenvalue-one restriction nor
subgroup cohomological dimension achieves this. The general conjecture
remains unresolved by this attempt.
\subsection{Sources}
\begin{itemize}
\item[S1]William M. Goldman, "Geometric Structures on Manifolds," AMS Graduate Studies in Mathematics 227 (2022), Conjecture 8.6.2; cf. H. Abels, G. Margulis, G. Soifer, arXiv:2011.12788 (Auslander conjecture for n < 7).
\url{https://www.math.umd.edu/~wmg/gstom.pdf}.
\textit{Formulation source.}
\item[E1]Herbert Abels, Grigory A. Margulis and Gregory A. Soifer,
\emph{The Auslander conjecture for dimension less than 7},
arXiv:2011.12788, 24 November 2020.
\url{https://arxiv.org/abs/2011.12788};
\url{https://arxiv.org/html/2011.12788v1}.
\textit{Primary abstract and introduction inspected 27 September 2026.
The dimension bound is $n<7$, not a theorem in arbitrary dimension.}
\item[E2]William M. Goldman, author draft of \emph{Geometric structures
on manifolds}, dated 11 December 2021, Chapters 8 and 15.
\url{https://www.math.umd.edu/~wmg/gstom.pdf}.
\textit{Author-hosted text inspected 27 September 2026. The date of
this accessible draft is distinguished from the 2022 book in [S1].}
\end{itemize}
\end{document}
